\exercisesection{3}

\begin{enumerate}
\def\labelenumi{\arabic{enumi}.}
\item
  对于一个整环成为因子分解整环，必要且充分条件是存在一个从 A - 0\(x \mapsto s(x)\) 到 \{N\} 的映射 \(s(xy) = s(x) + s(y)\)，使得 ，关\(s(x) = 0\)系  蕴含 x 在 A 中可逆；最后，对于 A - 0 中任\{\}意两个元素 x、y（它们互不整除），存在 A - 0 中\{\}的元素 \(ax + by = zt\)a\(s(z) < s(x)\)、b、z、t，使得 ，，且 t 与 x、y 互素。（若此条件满足，先证明 A 的每个非零元素都是极端元之积，再证明对每个极端元 p，Ap 是素理想；借助定理 1 (d) 得出结论）。
\item
  \begin{enumerate}
  \def\labelenumii{(\alph{enumii})}
  \tightlist
  \item
    设 A 是一个整环，它是右有向子环族  的并\((\mathbf{A}_{\lambda})\)。假设每个  都是因子分解整环\(A_{\lambda}\)，并且当  时，A＿\(\lambda \leqslant \mu\)bda 的每个极端元在\(A_{\lambda}\) A＿ 中\(A_{\mu}\)仍为极端元。证明 A 是因子分解整环，且其极端元集合是各 A＿ 的极端元集合之并（参见 § 2，习\(A_{\lambda}\)题 19 (e)）。
  \end{enumerate}
\end{enumerate}

\begin{enumerate}
\def\labelenumi{(\alph{enumi})}
\setcounter{enumi}{1}
\item
  由 (a) 推出：对于每个因子分解整环 \(\mathbf{A}\)，任意一族不定元上的多项式整环 \(\mathbf{A}[\mathbf{X}_{\lambda}]_{\lambda \in \mathbf{L}}\) 都是因子分解整环。
\item
  设 A 是一个因子分解整环，并且有限个不定元的每个形式幂级数整环 \(A[[X_{1},\ldots,X_{n}]]\) 都是因子分解整环。证明任意一族不定元的形式幂级数整环 \(A[[X_{\lambda}]]_{\lambda\in L}\)（《代数》第四章，§5，习题 1）也是因子分解整环（方法相同）。
\end{enumerate}

\begin{enumerate}
\def\labelenumi{\arabic{enumi}.}
\setcounter{enumi}{2}
\tightlist
\item
  \begin{enumerate}
  \def\labelenumii{(\alph{enumii})}
  \tightlist
  \item
    设 \(A = \bigoplus_{n \geqslant 0} A_n\) 是域 k 上具有正次数的分次代数；假设 \(A_0 = k\)，且 A 是 Krull 整环。证明 A 成为因子分解整环的必要且充分条件是：A 的每个高度为 l 的分次素理想 p 都具有 p = Aa 的形式，其中 a 是一个齐次元素（使用 §1 的习题 16）。证明 A 的每个 \(\neq 0\) 齐次元素都是齐次极端元之积。
  \end{enumerate}
\end{enumerate}

\begin{enumerate}
\def\labelenumi{(\alph{enumi})}
\setcounter{enumi}{1}
\tightlist
\item
  设\(\mathbf{A}\)  \(k\)是满足 (a) 条件的分次 k\(k'\)\(k\)\(\mathbf{A} \otimes_k k'\)\(\mathbf{A}\)\(\mathfrak{a}\)\(\mathbf{A}\)\(\mathfrak{a} \otimes_k k'\)\(\mathbf{A} \otimes_k k'\)\(\mathfrak{a}\)-代数。设  是 k 的扩域，并假设  是因子分解整环；证明  是因子分解整环。（若  的一个分次理想  使得  在  中是主理想，证明  是主理想；然后使用 (a)。）
\end{enumerate}

\begin{enumerate}
\def\labelenumi{\arabic{enumi}.}
\setcounter{enumi}{3}
\tightlist
\item
  \begin{enumerate}
  \def\labelenumii{(\alph{enumii})}
  \tightlist
  \item
    证明环 ，其\(\mathbf{A} = \mathbf{Q}[\mathbf{X},\mathbf{Y}] / p\)\(p\)中 p 是\(\mathbf{X}^2 +\mathbf{Y}^2 -1\)\(\mathbf{Q}[\mathbf{X},\mathbf{Y}]\) \(x\)\(\mathbf{X}\)\(\mathbf{A}\)\(x\)\(\mathbf{A}\)\(\mathbf{A}x\)\(\mathbf{A}\) 中由  生成的主理想，是一个非因子分解的 Krull 整环（若 x 是  在  中的像，证明 x 是  的极端元，但  不是  的极大理想）。
  \end{enumerate}
\end{enumerate}

\begin{enumerate}
\def\labelenumi{(\alph{enumi})}
\setcounter{enumi}{1}
\tightlist
\item
  证明环 \(\mathbf{A} \otimes_{\mathbf{Q}} \mathbf{Q}(i)\) 是因子分解整环（证明该环同构于 \(\mathbf{Q}(i) [X, Y]\) 除以主理想（ \(XY - 1\) ）所得的商；与习题 3 (b) 比较）。
\end{enumerate}

¶ 5. (a) 设 k 是一个 Noether 因子分解整环，B 是多项式环 \(k[X_{1},\ldots,X_{n}]\)，其中 \(n\geqslant3\)；设 \(g_{i}\;(0\leqslant i\leqslant r)\) 是 \(k[X_{3},\ldots,X_{n}]\) 中的元素

其中\(g_0\)  是极端元\(g = \mathbf{X}_1\mathbf{X}_2 - \sum_{i=0} g_i\mathbf{X}_1^i\)。令 \(\mathbf{A} = \mathbf{B}/g\mathbf{B}\)\(\mathbf{A}\)\(\mathbf{S}\)，\(\mathbf{A}\)\(\mathbf{X}_1\)\(\mathbf{A}\)\(\mathbf{S}^{-1}\mathbf{A}\)考虑商环 。证明  是因子分解整环。（令  为  中由 1 和  的像生成的乘法子集；对  应用第 4 节命题 3。）

\begin{enumerate}
\def\labelenumi{(\alph{enumi})}
\setcounter{enumi}{1}
\item
  \(k\)设 k 是特征\(\neq 2\)\(F\)  的域，F\(k[\mathbf{X}_1, \ldots, \mathbf{X}_n]\)\(n \geqslant 5\) 是\(k^n\) \(k[\mathbf{X}_1, \ldots, \mathbf{X}_n] / (\mathbf{F})\)\(G\)\(k[\mathbf{X}_1, \ldots, \mathbf{X}_n]\) 中的二次齐次多项式，其中 \(n \geqslant 3\)\(k\)，并且相应的 k＾n 上多项式函数是非退化二次型。证明整环  是因子分解整环。（首先证明：若  是  中的二次齐次多项式，且相应多项式函数是非退化二次型，则当  时 G 是极端元。然后在 k 代数闭时使用 (a) 证明命题；最后借助习题 3 (b) 推到一般情形。）
\item
  若 \(F = X_{1}X_{2} - X_{3}X_{4}\)，证明整环 \(k[X_{1}, X_{2}, X_{3}, X_{4}]/(F)\) 不是因子分解整环。（证明该环中各 \(X_{i}\) 的像都是极端元。）
\end{enumerate}

¶ 6. (a) 设 K 是一个特征  的代数闭域。确定环

\[
\mathrm{K} [ \mathrm{X}, \mathrm{Y}, \mathrm{Z} ] / (\mathrm{X} ^ {2} + \mathrm{Y} ^ {2} + \mathrm{Z} ^ {2}).
\]

\begin{enumerate}
\def\labelenumi{(\alph{enumi})}
\setcounter{enumi}{1}
\item
  设\(k\) k 是一个有\(a, b, c\)\(>0\)\(k\)\(A\)\(k[X, Y, Z] / (aX^2 + bY^2 + cZ^2)\)\(A\)序域，a、b、c 是 k \(k\)中的  元素，A 是环 。证明 A 是因子分解整环\(A\)。（使用习题 3 (b) 将其化为 k 是极大有序域的情形；然后利用 (a) 证明 A 的每个高度 1 的分次素理想都是主理想，并应用习题 3 (a)。）
\item
  设\(k\)  是一个有序域，a、b \(a, b\)是 k 中\(> 0\)的 \(k\) 元素\(B\)，B 是环
\end{enumerate}

\[
k [ \mathrm{X}, \mathrm{Y} ] / (a \mathrm{X} ^ {2} + b \mathrm{Y} ^ {2} + 1).
\]

证明 B 是因子分解整环（使用 (b) 和 § 1 的习题 17 (a)）。

\begin{enumerate}
\def\labelenumi{(\alph{enumi})}
\setcounter{enumi}{3}
\tightlist
\item
  证明整环 \(\mathbf{C} = \mathbf{Q}[\mathrm{X},\mathrm{Y}] / (\mathrm{X}^2 +2\mathrm{Y}^2 +1)\) 是因子分解整环，但整环 \(\mathbf{C}\otimes_{\mathbf{Q}}\mathbf{Q}(i)\)（其中 \(i^2 = -1\)）不是因子分解整环。
\end{enumerate}

¶ 7. 设 K 是一个 Noether 因子分解整环，F 是多项式\(K[X_{1},\ldots,X_{n}]\)环  中的极端元；假设当\(q(i)>0\)给\(1\leqslant i\leqslant n\)每个  赋予权 （\(X_{i}\)）时，F 是等重的且权为\textgreater q0。设 A 是分式域  的代数闭包  \(\Omega\)中由 K、（）以及多项式\(K(X_{1},\ldots,X_{n})\)  的一\(X_{i}\)个根\(1\leqslant i\leqslant n\) z 生成的整环，其中 \(Z^{c}-F\)c 是与 q 互素的整数。证明以下两种情形中 A 是因子分解整环：

\begin{enumerate}
\def\labelenumi{(\arabic{enumi})}
\item
  \(c \equiv 1 \pmod{q}\)；
\item
  每个有限生成的射影 K-模都是自由的（例如 K 是域、主理想整环或局部环时成立）。（第一种情形，考虑分式环 A[1/z]；证明它是因子分解整环{}，{并}应用第 4 节命题 3。第二种情形，取整数 d 使 ，并取  中满足  的 ；整环 B =\(cd \equiv 1 \pmod{q}\) A[\(z' \in \Omega\)z'] 由\(z = z'^{d}\)第一种情形是因{}子{}分解的，并且是自由 A-模。将 B 看作分次环，令  的权为 q，每个  的\(z'\)权为 ，并将问题化为证\(X_{i}\)明 A \(cdq(i)\)的两个齐次元素 u、v 的理想 Au ∩ Av 是主理想，使用 §1 的习题 16；最后考虑 B 中的理想 Bu ∩ Bv，并使用第一章，§3，第 6 节，命题 12。）特别地，如果 K 是域，且 a、b、\textgreater 是两两互素的三个整数 0，则整环
\end{enumerate}

\[
\mathbf {A} = \mathrm{K} [ \mathrm{X}, \mathrm{Y}, \mathrm{Z} ] / (\mathrm{Z} ^ {a} - \mathrm{X} ^ {b} - \mathrm{Y} ^ {c})
\]

是因子分解的。

¶ 8. 设 A 是一个整环，x、y、z 是 A 的三个非零元素，其中 x 是极端元且 Ax ∩ Ay = Axy。设 S 是由 \(x^{n}\)（\(n \geqslant 0\)）组成的乘法集，并令 B = S \(^{-1}\) A；考虑形式幂级数整环 A{[}{[}T{]}{]} 和 B{[}{[}T{]}{]}。

\begin{enumerate}
\def\labelenumi{(\alph{enumi})}
\tightlist
\item
    设 \(i, j, k\) 是三个满足 \(\geqslant 0\) 的整数，使得 \(ijk - ij - jk - ki \geqslant 0\) 且 \(z^i \in Ax^j + Ay^k\)。在 A{[}{[}T{]}{]} 中考虑元素 \(v = xy - z^{i-1}\mathrm{T}\)。证明存在一个整数 \(t > 0\) 和一个级数
\end{enumerate}

\[
v ^ {\prime} = y ^ {t} x ^ {- 1} + b _ {1} x ^ {- 2} \mathrm{T} + \dots + b _ {n - 1} x ^ {- n} \mathrm{T} ^ {n - 1} + \dots
\]

在 {}{}B{}{}[[T]]\(vv' \in A[[T]]\) 中，使得 \(b_{n}\)（递归确定 ，在 N 中按长度为 ij 的区间推进；在每个区间的内部，取

\[
b _ {n + 1} = b _ {n} z ^ {i - 1} x ^ {- 1};
\]

在每个区间内部，取\(ijk + ij - jk - ki \geqslant 0\) 

\begin{enumerate}
\def\labelenumi{(\alph{enumi})}
\setcounter{enumi}{1}
\item
  假设 \(z^{i-1} \notin Ax + Ay\)。证明在环 A[[T]] 中{不}{}{存}{}在常数项为 （k 是  的整数\(y^k\)）\(k\)且在 B\(>0\)[[T]] 中与 v 相伴的形式幂级{}{数}{}{（}计算 v 与 B[[T]] 中一个可逆元素的乘积的 T 系{数}{}）{}{}。
\item
  由 (a)、(b) 和习题 7 推出一个因子分解整环 A 的例子，使得形式幂级数整环 A[[T{}{]}{}{]} 不是因子分解的。（在上述记号中，并假定关于 x、y、z、i、j、k 的 (a)、(b) 条件均满足，证明  \(vv'\)不可能在 A[[T]] 中分解为极\(u_{h}\)端因\(1 \leqslant h \leqslant r\)子 {（}{}）{}{之}积；注意  是\(u_{h}\)常数项为 y 的幂的形式幂级数。然后考虑环 C = S'-1A[[T]]，其\^{}\{中\} {}{}S{}{'} 由常数项属于 S 的级数组成；B[[T]] 是 C{}{ }{的}{} Zariski 环的完备化；证明  属于 \(v' \in C\)C，且 v 和 \(u_{h}\) 在 C 中是极端元；最后由 (b) 得到矛盾。）
\item
  由 (c) 推出一个 Noether 因子分解局部整环 A 的例子，其完备化 \(\hat{A}\) 是非因子分解的 Krull 整环。
\end{enumerate}

¶ 9. (a) 设 A 是一个 Noether 整环，使得对 A 的每个极大理想 m，\(A_{m}\) 都是离散赋值环。若 B 是形式幂级数环 \(A[[X_{1},\ldots,X_{n}]]\)，证明对 B 的每个极大理想 n，整环 \(B_{n}\) 都是因子分解整环（考虑其完备化，使用第 9 节命题 8）。推出 B 的每个除性理想都是投射 B-模。

\begin{enumerate}
\def\labelenumi{(\alph{enumi})}
\setcounter{enumi}{1}
\item
  设 C 是一个 Noether 环，使得每个有限生成投射 C-模都是自由的；证明形式幂级数环 C{[}{[}X{]}{]} 具有相同性质（参见第二章，§ 3，第 2 节，命题 5）。
\item
  由 (a) 和 (b) 推出：如果 A 是主理想整环，则形式幂级数整环 \(A[[X_{1},\ldots,X_{n}]]\) 是因子分解的。
\end{enumerate}

\begin{enumerate}
\def\labelenumi{\arabic{enumi}.}
\setcounter{enumi}{9}
\tightlist
\item
  \begin{enumerate}
  \def\labelenumii{(\alph{enumii})}
  \tightlist
  \item
    Prüfer（§ 2，习题 12）因子分解整环是主理想整环。
  \end{enumerate}
\end{enumerate}

\begin{enumerate}
\def\labelenumi{(\alph{enumi})}
\setcounter{enumi}{1}
\tightlist
\item
  伪 Bézout（§ 1，习题 21）Krull 整环是因子分解的。
\end{enumerate}

\begin{enumerate}
\def\labelenumi{\arabic{enumi}.}
\setcounter{enumi}{10}
\item
  设 K 是一个域，A 是多项式整环 \(K[X,Y]\)，它是因子分解的；若 L 是域 \(K(X^{2},Y/X) \subset K(X,Y)\)，证明整环 \(A \cap L\) 不是因子分解的。
\item
  使用第三章，§ 2，第 8 节，定理 1 的推论 3 证明第 8 节命题 5。
\item
  将第 8 节命题 7 的推论推广到完备 Hausdorff 局部环 A 不是整环的情形（使用《代数》第八章，§6，习题 6 (b)）。
\item
  设 A 是一个完备 Noether 局部环，其剩余域的特\textgreater{}征为 。在形式幂级数环 A[{}{}[{}{}T]] 中，对每\(\omega_{n} = (1 - T)^{p^{n}}\)个整\(\gamma_{n} = 1 - \omega_{n}\)数 ，考虑元素\textgreater{}  和 。\(\gamma_{n}\)证明  除了一个符号因子外是一个特异多项式（第 \(A_{n} = A[[T]]/(\gamma_{n})\)8 节）；推出  可识别为 A 上群
\end{enumerate}

\(G_{n} = Z/p^{n}Z\)。证明主理想 \((\gamma_{n})\) 的交为 0；推出 A{[}{[}T{]}{]} 可识别为逆极限 \(\lim A_{n}\)。

¶ 15. 设 K 是一个关于赋值 v 完备的域，A 是该赋值的环，k 是其剩余域，P 是 A[X₁, , Xₙ] 中总次{}数为\ldots d{} 的多项式，满足以下性质：存在 K' 的一个代数扩张，使得在 K'[X₁, , Xₙ] 中{} P\ldots {是}总次数为 1 的多项式之积。进一步假设 A[X₁, , Xₙ] 中存在两个多项式 Q、R，{}使 \ldots {}的总次数为 s，且含有单项式 ，其中 （ 表示典\(aX_1^{s}\)范同态 A → k），R 的次数 ≤d - s，且  . （第三章，§ 4\(\overline{\mathbf{P}} = \overline{\mathbf{Q}}\) 的\(\overline{\mathbf{R}}\)记号）。证明于是存在 A[X₁, , Xₙ] 中的两{个}多\ldots 项式{} Q₀、R₀，其次数分别为 s、d - s，使得 P = Q₀R\(\overline{\mathbf{Q}} = \overline{\mathbf{Q}}_0\)₀\(\overline{\mathbf{R}} = \overline{\mathbf{R}}_0\)，，，且 Q₀ 含有单项\(a_0X_1^{s}\)式 ，其中 。（将 P、Q、R 看作关于 X₁ 的多项式，其系数在环 B 中；B 是 A[X₂{}, \ldots, {}Xₙ] 关于通过第六章，§ 10，第 1 节，命题 2 的方法将 v 延拓所得赋值的完备化；然后应用 Hensel 引理；最后使用 P 的初始假设。）

\begin{enumerate}
\def\labelenumi{\arabic{enumi}.}
\setcounter{enumi}{15}
\item
  设 B 是一个剩余域为有限域且不是素域的离散赋值环；设  是\(k_{0}\) k 的素子域，A 是 B 的子环，由剩余域中类属于  的元素组成。设  是 B 的\(k_{0}\)一个\(\pi\)一致化元， 是 \((\theta_{i})_{1\leq i\leq m}\)B 的一组可逆元素，使得  的模  的类\(\bar{\theta}_{i}\)  \(\pi\)构成\(\theta_{i}\) k＾* 模  的代表系\(k^{*}\)。证\(k_{0}^{*}\)明元素  和元素\(p_{i}=\theta_{i}\pi\)  在 A \(\pi\)中都是极端元，且 A 中每个元素都是一个可逆元素与  及  的幂的\(p_{i}\)乘\(\pi\)积，尽管 A 不是整闭的。
\item
  \begin{enumerate}
  \def\labelenumii{(\alph{enumii})}
  \tightlist
  \item
  设 A 是一个整环；证明在环  \(A[X_{ij}]\)中，其\((\mathbf{X}_{ij})\)中  是由 \(n^{2}\)n² 个不\((1 \leqslant i \leqslant n, 1 \leqslant j \leqslant n)\)定元组成的\(\det(\mathbf{X}_{ij})\)族（，），元素  是极端元。（将其化为 A 是域的情形；注\(\det(\mathbf{X}_{ij})\)意  的因子必为齐次多项式，然后对 n 作归纳。）
  \end{enumerate}
\end{enumerate}

\begin{enumerate}
\def\labelenumi{(\alph{enumi})}
\setcounter{enumi}{1}
\tightlist
\item
  \(\mathbf{K}\)设  是无限域， \(\mathbf{F}\)是  中的\(\mathbf{K}[\mathbf{Y}_1,\dots ,\mathbf{Y}_m]\)多项式，也写作 ；\(\mathrm{F(Y)}\)对于 K 上任意一\(s = (\alpha_{ij})\)个 m\(m\) 阶方阵 ，记 \(F(s.Y)\) 为如下多项式 F：将元素  \(\sum_{j=1}^{m}\alpha_{ij}Y_{j}\)代入每个 。证明若 \(Y_{i}\)F 是极端元，则对于每个\(F(s.Y)\)可逆矩阵 s， 也是极端元。若存在整数 \(k\geqslant0\) 使得
\end{enumerate}

\[
\mathbf {F} (s. \mathbf {Y}) = (\det (s)) ^ {k} \mathbf {F} (\mathbf {Y})
\]

对于每个可逆矩阵\(s\) \(\mathbf{F}\)s， 必然关于每个  齐次\(\mathbf{Y}_j\)；

此外，若 ，其\(\mathbf{F} = \mathbf{GH}\)中  \(\mathbf{G}\)和  \(\mathbf{H}\)是  中的两个多项式，则\(\mathbf{K}[\mathbf{Y}_1,\dots ,\mathbf{Y}_m]\)存在两个整数 p、q，使得 p\(p\) +\(q\) q = k\(p + q = k\) 且

\[
\mathbf {G} (s. \mathbf {Y}) = (\det (s)) ^ {p} \mathbf {G} (\mathbf {Y}), \quad \mathbf {H} (s. \mathbf {Y}) = (\det (s)) ^ {q} \mathbf {H} (\mathbf {Y})
\]

对于每个可逆矩阵 s（使用 \(s\)(a)）。

\begin{enumerate}
\def\labelenumi{(\alph{enumi})}
\setcounter{enumi}{2}
\tightlist
\item
  设 A 是一个不为 0 的环；考虑多项式环 ，其中  是 （分别为 \(\mathbf{A}[\mathbf{X}_{ij}]\)）个不定元，\(\mathbf{X}_{ij}\)满足\(n(n + 1)/2\) （分别\(2n(2n - 1)/2\)为 ）；设 （分别为 \(1 \leqslant i \leqslant j \leqslant n\)）是  \(1 \leqslant i < j \leqslant 2n\)上的 n\(U = (\xi_{ij})\) 阶（分\(V = (\eta_{ij})\)别为 2n 阶）方阵，满足：当 \(n\) 时 \(2n\)，且当 \(\mathbf{A}[\mathbf{X}_{ij}]\)i > j \(\xi_{ij} = \mathbf{X}_{ij}\)时 \(1 \leqslant i \leqslant j \leqslant n\)（分\(\xi_{ij} = \mathbf{X}_{ji}\)别地，\(i > j\)当  时\(\eta_{ii} = 0\) ，\(1 \leqslant i \leqslant 2n\)当\(\eta_{ij} = \mathbf{X}_{ij}\)  时\(1 \leqslant i < j \leqslant 2n\) \(\eta_{ij} = -\mathbf{X}_{ji}\)，当 \(i > j\)i > j \(\det(U)\)时 ）。\(\operatorname{Pf}(V)\)证明 （分别为 ）是  中的\(\mathbf{A}[\mathbf{X}_{ij}]\)极端元（仿照 (b)，考虑  \(\det(s, U, {}^ts)\)和 ）\(\operatorname{Pf}(s, V, {}^ts)\)。
\end{enumerate}

\begin{enumerate}
\def\labelenumi{\arabic{enumi}.}
\setcounter{enumi}{17}
\tightlist
\item
  设 K 是一个域，f = g/h 是二元不定元的有理函数域 \(\mathrm{K}(\mathrm{U}, \mathrm{V})\) 中的元素，其中 g、h 是 \(\mathrm{K}[\mathrm{U}, \mathrm{V}]\) 中互素的两个多项式。证明在有理函数域
\end{enumerate}

\[
\mathrm{K} \left(\mathrm{X} _ {1}, \mathrm{Y} _ {1}, \dots , \mathrm{X} _ {n}, \mathrm{Y} _ {n}\right)
\]

在含有 \(2n\) 个不定元的有理函数域中，行列式 \(\operatorname{det}(f(\mathbf{X}_i, \mathbf{Y}_j))\) 等于：

\[
\left(\prod_ {i, j} h \left(\mathrm{X} _ {i}, \mathrm{Y} _ {j}\right)\right) ^ {- 1} \mathrm{V} \left(\mathrm{X} _ {1}, \dots , \mathrm{X} _ {n}\right) \mathrm{V} \left(\mathrm{Y} _ {1}, \dots , \mathrm{Y} _ {n}\right) \mathrm{F} \left(\mathrm{X} _ {1}, \mathrm{Y} _ {1}, \dots , \mathrm{X} _ {n}, \mathrm{Y} _ {n}\right)
\]

其中 F 是  中的多项式， \(K[X_{1}, Y_{1}, \ldots, X_{n}, Y_{n}]\)是 V\(V(X_{1}, \ldots, X_{n})\)andermonde 行列式（《代数》第三章，§ 6，第 4 节）。考虑 f = 1/(U + V) 的特殊情形（\(f = 1/(U + V)\)“Cauchy 恒等式”）。

\begin{enumerate}
\def\labelenumi{\arabic{enumi}.}
\setcounter{enumi}{18}
\tightlist
\item
  \(U\)若 U 是 n 阶方阵，\(n\) \(\Delta\)是其行列式， 是\(\Delta_p\) U 的第 p 个外\(p\)幂的行列式（《代数\(U\)》第三章，§ 6，第 3 节），证明：
\end{enumerate}

\[
\Delta_ {p} = \Delta^ {\binom {n - 1} {p - 1}}
\]

（使用《代数》第三章，§ 6 的习题 11 和上面的习题 17 (a)。）

\begin{enumerate}
\def\labelenumi{\arabic{enumi}.}
\setcounter{enumi}{19}
\tightlist
\item
  \(\mathbf{A}\)设  是一个因子分解整\(f = \sum_{k=0}^{n} a_k \mathbf{X}^k\)环， 是  \(\mathbf{A}[\mathbf{X}]\)中的多项式；假设 A 中存在一个极端元\(p\) p\(\mathbf{A}\)，使得：
\end{enumerate}

\begin{enumerate}
\def\labelenumi{(\arabic{enumi})}
\item
  存在一个索引 ，使 a\(k \leqslant n\)＿k 不被 \(a_{k}\)p 整除，但当 i <\(p\) k \(a_{i}\)时 a＿i 被 p\(p\) 整除\(i < k\)；
\item
 \(a_0\) a＿0 被  整除，\(\pmb{p}\)但不被  整除\(p^2\)。
\end{enumerate}

证明在这些条件下，f 在  中的不可约因子之一的次数 \(f\)（在\(\mathbf{A}[\mathbf{X}]\)  中论证\(\geqslant k\)）。考虑 \((\mathbf{A} / p\mathbf{A})[\mathbf{X}]\)k = n 的特殊情形（“Eis\(k = n\)enstein 不可约性判别法”）。

\begin{enumerate}
\def\labelenumi{\arabic{enumi}.}
\setcounter{enumi}{20}
\tightlist
\item
  证明在\(\mathbf{Z}[\mathbf{X}]\)  中下列多项式不可约：，其中 \(\mathbf{X}^n - a\)a 的一个素因子指数为 1；\(a\)（将  替换\(\mathbf{X}^{2^k} + 1\)为 ）；\(\mathbf{X}\)；。\(\mathbf{X} + 1\)（\(\mathbf{X}^4 + 3\mathbf{X}^3 + 3\mathbf{X}^2 - 5\)使\(5\mathbf{X}^4 - 6\mathbf{X}^3 - a\mathbf{X}^2 - 4\mathbf{X} + 2\)用习题 20。）
\end{enumerate}

¶ 22. (a) 设 k 是一个有序域，A 是多项式环 k{[}X, Y, Z{]} 除以主理想 \((\mathbf{X}^{2} + \mathbf{Y}^{2} + \mathbf{Z}^{2} - 1)\) 所得的商；设 x、y、z 表示 X、Y、Z 在 A 中的典范像。证明 A 是因子分解整环（考虑分式环 \(A_{z-1}\)（第二章，§ 5，第 1 节的记号），使用第 4 节命题 3）。

\begin{enumerate}
\def\labelenumi{(\alph{enumi})}
\setcounter{enumi}{1}
\tightlist
\item
  \((e_i)_{1 \leq i \leq 3}\)设  是  的典\(\mathbf{A}^3\)范基，\(\mathbf{M}\) 是  除以由\(\mathbf{A}^3\)  生成的单生成子 \(\mathbf{N}\)A-模 \(xe_1 + ye_2 + ze_3\) 所得的\(\mathbf{M}\)商；证明  是投射 A-模（在  中为\(\mathbf{N}\)  \(\mathbf{A}^3\)构造一个补子模\(k = \mathbf{R}\)）。* (\(\mathbf{M}\)c) 证明若 ，\(\mathbf{M}\)A-模  不是自由的（将  识别为单位球面切向量丛的连续截面模的子模，并使用球面上不存在在每一点都不为 0 的连续切向\(\neq 0\)量场这一事实）。 *
\end{enumerate}

\begin{enumerate}
\def\labelenumi{\arabic{enumi}.}
\setcounter{enumi}{22}
\tightlist
\item
  设 B 是 Krull 整环，E 是其分式域，\(\Delta\) 是 E 的一个导子，满足 \(\Delta(\mathbf{B}) \subset \mathbf{B}\)，K 是 E 中 \(\Delta\) 的核所成的子域，A 是 Krull 整环 \(B \cap K\)；假设 K 的特征为 p \textgreater{} 0；则 \(E^{p} \subset K\) 且 \(B^{p} \subset A\)，因此 B 是 A 在 E 中的整闭包，且典范同态 \(\vec{i}: C(A) \to C(B)\) 有定义（\(\S 1\)，第 10 节）。记 U 为 B 的可逆元素群。
\end{enumerate}

\begin{enumerate}
\def\labelenumi{(\alph{enumi})}
\item  若  满足  是  中某个除子的典范像，证明 （注意，对  的每个高度为 1 的素理想 ，存在 ，使得 ，并注意  在  下不变）。令  为  的加法子群，由属于  的 （“对数导数”）组成（对于  或 ，这两者等价，且 ），并令  为  的子群，由 （其中 ）组成。证明：对于每个 ，若 d 的像是  中的主除子，则对所有满足  的 b， 在模  下的类只取决于 d 在  中的类，并导出从  到  的典范单射同态 。
  \(b \in \mathbf{E}\)若  满足\(\mathrm{div}_{\mathbf{B}}(b)\)  是  中某个除子的典范像\(\mathbf{D}(\mathbf{A})\)，证明 \(\Delta b / b \in \mathbf{B}\)（注意，对  的每个高度\(\mathfrak{P}\)为 \(\mathbf{B}\)1 的素理想 ，存\(b' \in \mathbf{K}\)在 ，使\(v_{\mathfrak{P}}(b) = v_{\mathfrak{P}}(b')\)得 ，并注意\(\mathbf{B}_{\mathfrak{P}}\)  在  下不\(\Delta\)变）。\(\mathbf{L}\)令  为  的加法子\(\mathbf{B}\)群，由属于  \(\Delta b / b\)的 （“对数导数”）组成（对于 \(\mathbf{B}\) 或\(b \in \mathbf{E}\) ，\(b \in \mathbf{B}\)这两者等价，且 ），并令\(b \neq 0\)  为\(\mathbf{L}'\)  的子群，由\(\mathbf{L}\) （其中 ）组\(\Delta u / u\)成。证\(u \in \mathbf{U}\)明：对于每个 ，若 d\(d \in \mathbf{D}(\mathbf{A})\) 的像是  中的主\(d\)除子，则对所有满足\(\mathbf{D}(\mathbf{B})\)  的 b， \(\mathbf{L}'\)在\(\Delta b / b\)模  \(b\)下的类只\(i(d) = \mathrm{div}_{\mathbf{B}}(b)\)取决于 d 在  中的\(d\)类\(\mathbf{C}(\mathbf{A})\)，并导出从  到  的典范单射\(\phi\)同态\(\operatorname{Ker}(\bar{i})\) \(\mathbf{L} / \mathbf{L}'\)。
\item
  证明：若\(\Delta(B)\)  不包含于 B 的任何\(B\)高度为 1 \([E: K] = p\)的\(\phi\)素理想，且 ，则  是双射。（化为\(b \in E\)证明：若 \(\Delta b / b \in B\) 满\(\mathfrak{P}\)足 ，且  是\(B\) B 的高度为 \(v_{\mathfrak{P}}(b)\)1 的素理想，使\(p\)得  \(e(\mathfrak{P} / \mathfrak{p}) = 1\)不是 \(\mathfrak{p} = \mathfrak{P} \cap A\)p 的倍数，则 ，其中 ；为此，由假\(t\)设推出：若 t\(B_{\mathfrak{P}}\) 是 \(\Delta t / t \in B_{\mathfrak{P}}\) 的一致\(\mathfrak{P}B_{\mathfrak{P}}\)化元，则 ，从\(\Delta\)而  在\(\Delta\)  下不变，因而  通过取商在剩余域\(\overline{\Delta}\)  上定义一个导\(k = B_{\mathfrak{P}} / \mathfrak{P}B_{\mathfrak{P}}\)子；证明 \(\overline{\Delta} \neq 0\)，并推出 。\(f(\mathfrak{P} / \mathfrak{p}) = p\))
\item
  \(k\)设 k 是特征 2 的域，B 是多项\(k[\mathbf{X}, \mathbf{Y}, \mathbf{Z}]\)式环\(\Delta\) ， 是  \(\mathbf{E} = k(\mathbf{X}, \mathbf{Y}, \mathbf{Z})\)的导子\(\Delta(\mathbf{X}) = \mathbf{Y}^4\)，\(\Delta(\mathbf{Y}) = \mathbf{X}^2\)满\(\Delta(\mathbf{Z}) = \mathbf{XYZ}\)足 、、；\(\mathbf{K}\)域  是 \(\Delta\) 的核，\([\mathbf{E} : \mathbf{K}] = 4\)并满足\(\Delta(\mathbf{Z}) / \mathbf{Z} \in \mathbf{B}\) 。于是 \(\mathrm{div}_{\mathbf{B}}(\mathbf{Z})\)，但证明  不是  中某个除子\(\mathrm{D}(\mathbf{A})\)的典范像。（反证，假设对于 、，\(e(\mathfrak{P} / \mathfrak{p}) = 1\)有\(\mathfrak{P} = \mathbf{B}\mathbf{Z}\) \(\mathfrak{p} = \mathfrak{P} \cap \mathbf{A}\)；那么  中会有 \(\mathbf{A}\) 的一致化元\(\mathrm{B}_{\mathfrak{P}}\)，它必为  的形式\(a(\mathbf{X}, \mathbf{Y}, \mathbf{Z})\mathbf{Z}\)，其\(b = a(\mathbf{X}, \mathbf{Y}, 0) \neq 0\)中 ；推出\(\Delta b / b = -\mathbf{X}\mathbf{Y}\) ，并通过计算  得到矛盾。\(\Delta(-\mathbf{X}\mathbf{Y})\)）
\end{enumerate}

\begin{enumerate}
\def\labelenumi{\arabic{enumi}.}
\setcounter{enumi}{23}
\tightlist
\item
  \begin{enumerate}
  \def\labelenumii{(\alph{enumii})}
  \tightlist
  \item
    设 E 是特征 2 的域，\(\Delta\) 是 E 的导子，K 是 E 中  的核所成的\(\Delta\)子域；假设 \([E:K]=2\)。证明 ，\(\Delta^{2}=a\Delta\)其中 \(a\in K\)，并证明对于 E 中\(t\in E\)元素 t，要写成\(\Delta x/x\)  的形式，必要且充分条件是 \(\Delta t=at+t^{2}\)。
  \end{enumerate}
\end{enumerate}

\begin{enumerate}
\def\labelenumi{(\alph{enumi})}
\setcounter{enumi}{1}
\tightlist
\item
  设 B 是特征 2 的因子分解局部整环，m 是其极大理想，E 是其分式域， 是 E\(\Delta\) 的导子；假设 E 的子域 K（ 的核）满足 ；此\(\Delta\)外，假设 {}m{ }中存在两个元素 x、y，使  和  生成 B 的理想 \(\Delta x\)q，\(\Delta y\)该理想由  生成。证明：若  \(\Delta(B)\)属于 q，则存在\(t = \Delta z / z\) B 的一个可逆元素 u，使 （写成 ，其中 r\(t = \Delta u / u\)、s \(t = r \Delta x + s \Delta y\)属于 B，并使用 (a)）。
\end{enumerate}

¶ 25. 设 k 是特征 2 的域，B 是形式幂级数环 ，E 是\(k[[X,Y]]\)其分式域， 是由 、 定义的\(\Delta\) E 的 k-导子（i、j \(\Delta(X)=Y^{2j}\)是满\(\Delta(Y)=X^{2i}\)足  的整数\(\geqslant0\)）；B 中由满足  的  组成的子\(x\in B\)环 A，是\(\Delta x=0\)  中的形式幂级数环，其中以  代\(k[[X,Y,Z]]\)替 X\(X^{2}\)、以  代替 Y、\(Y^{2}\)以  代替\(X^{2i+1}+Y^{2j+1}\) Z。

\begin{enumerate}
\def\labelenumi{(\alph{enumi})}
\item
  证明 C(A) 包含一个维数为  的 k 上向量\(N(i,j)\)空间，该维数等于满足 、 且  的整\((a,b)\)数有序对\(0 \leqslant a < i, 0 \leqslant b < j\)  \((2j + 1)a + (2i + 1)b \geqslant 2ij\)的个数。（使用习题 23 (b) 和习题 24 (a)；注意 L 的元素\(F \in B\)是 B \(\Delta F = F^{2}\)中满足  的形式幂级数；\(2j + 1\)给 X 赋予权\(2i + 1\) ，给 Y 赋予权 ，将 F 分解为等重多项式\(L_{q}\)的无穷和；若  是 L 的子\(F \in L\)群，由其等重分量的权  的  \(\geqslant q\)组\(L/L'\)成，则  同构于群  的直和，其中\(C_{q}/C_{q+1}\) ；对 \(C_{q} = L_{q}/(L' \cap L_{q})\)，使用习题 24 (b\(q \geqslant 4ij\)) 计算这些群。）
\item
证明 A 的理想 \(AX^{2}\) 是素理想；若 \(A' = A[X^{-2}]\)，推出 C(A') 与 C(A) 同构（第 4 节命题 3）。证明 A' 是 Dedekind 整环（考虑环 \(B' = B[X^{-2}]\)，它整于 A' 且是主理想整环，并使用第五章，§ 2，第 4 节，定理 3）。推出一个 Dedekind 整环的例子，其理想类群是无限的。
\end{enumerate}

\begin{enumerate}
\def\labelenumi{\arabic{enumi}.}
\setcounter{enumi}{25}
\tightlist
\item
  设 A 是一个因子分解整环，K 是其分式域。要\(i\)使\(1 \leqslant i \leqslant r\) B {=}\(_{1}\) \ldots[\(_{n}\){}X  , , X  \(\sum_{i=1}^{r} \text{Bf}_i\)] 中的元素 （）满足理想  等于 B，必要且充分\(v_{i}\)条\(1 \leqslant i \leqslant r\)件是\(\mathbf{K}[\mathbf{X}_1, \ldots, \mathbf{X}_n]\)存在  \(\sum_{i=1}^{r} v_i f_i = 1\)中的多项式 （），使得 ，并且满足以下附加条\(v_i = w_i / d\)件：若\(d \in \mathbf{A}\)写 ，其中 \(w_i\)，多项式\(\mathbf{A}[\mathbf{X}_1, \ldots, \mathbf{X}_n]\)  属于 ，且所有  的系数集合的最\(w_i\)大\(1 \leqslant i \leqslant r\)公因子（）等于 1，则对于 A 的\(p\)每个整除 \(d\)d 的极端元 p， 在环  中\(f_i\)的类所生成\((\mathrm{A} / \mathrm{A}p)[\mathrm{X}_1, \ldots, \mathrm{X}_n]\)的理想就是整个环。
\end{enumerate}

¶ 27. (a) 设 K 是一个域，B 是在四个不定元 K(U, V, X, Y) 的有理函数域的代数闭包中由多项式环 K{[}U, V, X, Y{]} 和多项式

\[
\mathrm{F} = \mathrm{Z} ^ {7} - \mathrm{U} ^ {5} \mathrm{X} ^ {2} - \mathrm{V} ^ {4} \mathrm{Y} ^ {3}.
\]

证明 \(\mathbf{B}\) 是因子分解整环（参见习题 7）。

\begin{enumerate}
\def\labelenumi{(\alph{enumi})}
\setcounter{enumi}{1}
\tightlist
\item
  设 p 是 A 中由 X、Y、U、V 和 z 生成的（素）理想，并记 \(C = A_{p}[[T]]\)；证明 C 是一个 Noether 局部整环，不是因子分解的，但其伴随分次整环 \(\text{gr}(C)\) 是因子分解的（使用习题 8）。
\end{enumerate}
% semantic-fragments: legacy-input-seam
\relax{}
